\preliminar{Abstract}

\begin{otherlanguage}{english}
% El entorno abre grupo: la marca decimal inglesa no escapa de este apartado.
\sisetup{output-decimal-marker={.}}

\textit{Diffusion Policy} generates robotic manipulation actions through a diffusion
process conditioned on the observation. Its visuomotor variant chains a visual encoder and
a generative network, so the image representation bounds the information available for
control. Subsequent research has mostly modified the generative block, whereas the choice
of encoder and of its training strategy remains unsettled when only a small set of
demonstrations is available.

This work compares five visual blocks integrated as the observation encoder of a single
\textit{Diffusion Policy} on the Push-T task: a ResNet-18 trained from scratch, the same
network pretrained on ImageNet in frozen and fine-tuned versions, and two frozen vision
transformers, DINOv2 ViT-S/14 and CLIP ViT-B/16. All of them output a
\num{512}-component vector and share the generative network, the data, and the evaluation
protocol.
Training uses \num{90} of the \num{206} teleoperated demonstrations, following the
reference implementation, and budgets ranging from \num{200} to
\num{500} epochs, with early stopping. A set of \num{50} episodes generated with held-out
seeds selects the checkpoint of each variant. Performance is then measured once on
\num{200} episodes from a disjoint seed block, under a preregistered protocol, together
with the compute time consumed. Each visual block is trained exactly once, a design
premise adopted in order to cover five configurations within the available compute
budget; the unit of inference is therefore the trained artefact, not the training
strategy.

Each condition is solved with two realisations of the diffusion noise, so that the estimand
integrates both the initial condition and the stochastic sampling of the policy. On that
independent block, the baseline trained from scratch reaches \num{0.887}. The four
pretrained variants score \num{0.642}, \num{0.580}, \num{0.571}, and \num{0.490}. V0
outperforms all of them after the ten pairwise tests are corrected with Holm's procedure,
and this holds regardless of which test is used. Among the pretrained variants, by contrast,
only the weakest separates from the rest under both tests, and the two intermediate ones are
indistinguishable from each other. Each V0 epoch takes \SI{1.6}{\minute}, compared with
\SI{2.3}{\minute} to \SI{14.7}{\minute} for the other variants. The five runs total
\SI{237.1}{\hour}.

A separately preregistered contrast compares that baseline with the checkpoint released by
the authors of the reference implementation. The released model scores \num{0.846}, that is
\num{0.041} lower. With an equivalence margin of \num{0.05} fixed in advance, the data
support the claim that the reproduction does not underperform the reference artefact,
without establishing that it exceeds it by a practically relevant amount. That contrast also
revealed that changing only the noise seed shifts the mean score by \num{0.030}, a magnitude
comparable to the differences separating the pretrained variants from one another, and this
measurement is what prompted duplicating the realisations throughout the final test.

The study concludes that, among the five artefacts evaluated and over the conditions
tested, none of the pretrained visual blocks improves the baseline either in performance
or in cost per epoch. The advantage must be ascribed to the complete visual block
(initial weights, resolution, cropping, and spatial aggregation), which varies jointly,
rather than to pretraining in isolation. With a single training run per variant, the
result does not license extrapolation to the expected behaviour of those strategies upon
retraining.

\vspace{1em}
\noindent\textbf{Keywords:} imitation learning, diffusion models, visual encoder, transfer
learning, robotic manipulation.
\end{otherlanguage}
